\documentclass{name}
\usepackage{lipsum}
\usepackage{multicol}
\usepackage{listings}
\usepackage{longtable}
\usepackage{amsmath}
\usepackage{mathtools}
\usepackage{amssymb}
\usepackage{booktabs}
\usepackage{array}
\usepackage{caption}
\usepackage{subcaption}
\usepackage{float}
\usepackage{setspace}
\usepackage{stmaryrd}
\usepackage[backend=bibtex,sorting=none]{biblatex}
\addbibresource{bibliographie.bib}
\usepackage{csquotes}
\usepackage[english]{babel}
\usepackage{fancyhdr}
\usepackage{lmodern}
\usepackage{xcolor}
\usepackage{algorithm}
\usepackage{algorithmic}
\usepackage{awesomebox}
\usepackage{todonotes}
\usepackage{upquote}

\newtheorem{theorem}{Theorem}
\newtheorem{remark}{Remark}
\newtheorem{definition}{Definition}

\newcounter{pagezero}

\begin{document}

%----------- Informations du rapport ---------
\UE{S8 Machine Learning}
\sujet{Kaggle Project: Textile Classification \& AI Safety}
\titre{TILDA Textile Classifier and AI Safety}

\eleves{Marin \textsc{Decanini}\\Nolan \textsc{Carlisi}}

\enseignant{Pierre \textsc{Weiss}\\Université de Toulouse\\ Senior CNRS researcher at IRIT}

%----------- Initialisation -------------------
\fairemarges %Afficher les marges
\fairepagedegarde %Créer la page de garde

\newpage
\setcounter{pagezero}{0}

\begin{center}
	\begin{minipage}{13cm}
	    \begin{center}
	        \section*{Abstract}
	    \end{center}
	    This report presents our implementation and experimental analysis for the S8 Machine Learning Kaggle project on the TILDA textile dataset. 
	    We focus on two core components: first, training state-of-the-art Convolutional Neural Networks (CNNs) for classification; and second, analyzing machine learning security and AI safety under biased training distributions. 
	    Specifically, we implement LeNet and ResNet architectures, evaluate their performance using mini-batch SGD, and perform systematic hyperparameter tuning. 
	    Furthermore, we investigate model vulnerability to spurious correlations by constructing a biased dataset with a noise-like channel correlated with labels, measuring both classification accuracy and fairness using the Disparate Impact (DI) metric.
	\end{minipage}
\end{center}

\par 
\setcounter{page}{\value{pagezero}}

\newpage
\setcounter{page}{\value{pagezero}}
\tableofcontents

\newpage

\section{Introduction}
% Introduce the TILDA dataset and the goals of the Kaggle project.
This project focuses on image classification using the TILDA textile texture dataset, with the dual goals of:
\begin{enumerate}
    \item Achieving high classification accuracy on the test dataset using state-of-the-art Convolutional Neural Network (CNN) models.
    \item Studying machine learning vulnerability and AI safety, specifically focusing on how models can be biased towards spurious correlations when trained on biased datasets.
\end{enumerate}
In the following sections, we describe our image pre-processing, architecture choices, training optimization setups, and parameter tuning, before diving into the mathematical construction and empirical evaluation of model bias.

\section{Train State-of-the-Art CNN Models}

\subsection{Image Pre-processing and Data Augmentation}
% Question 2.1
To prepare the dataset for deep learning models, we implement a robust image pre-processing pipeline. 
\begin{itemize}
    \item \textbf{Image Resizing:} All input images from the TILDA dataset are resized to a uniform input shape (e.g., $128 \times 128$ or $224 \times 224$ pixels) to fit the requirements of the selected network architectures.
    \item \textbf{Normalization:} Channel-wise normalization is applied to standardize pixel values across the dataset.
    \item \textbf{Data Augmentation:} To improve model generalization and mitigate overfitting, we apply randomized transformations during training, such as random horizontal flips, vertical flips, random rotations, and mild color jittering.
\end{itemize}
The exact configuration of these pre-processing transforms is detailed in the experimental code.

\subsection{CNN Architectures Definition and Rationale}
% Question 2.2
We implement and compare two primary architectures:
\begin{enumerate}
    \item \textbf{LeNet-5} \cite{lecun1998gradient}: A classical architecture comprising convolutional layers, average/max pooling, and fully-connected layers. This serves as our baseline.
    \item \textbf{ResNet} \cite{he2016deep}: A modern architecture utilizing skip connections (residual blocks) to allow training of deeper networks by mitigating the vanishing gradient problem.
\end{enumerate}

\subsubsection{Architecture Details}
% Rationale and precise definition: number of layers, types, activation functions, and what each layer does.
\begin{table}[H]
\centering
\caption{Detailed Architectural Layouts for Chosen Models}
\label{tab:architectures}
\begin{tabular}{llll}
\toprule
Model & Layer Type & Specifications / Hyperparameters & Activation \\
\midrule
\textbf{LeNet-5} & Conv2d & In: 1/3, Out: 6, Kernel: $5 \times 5$, Stride: 1 & ReLU \\
 & MaxPool2d & Kernel: $2 \times 2$, Stride: 2 & - \\
 & Conv2d & In: 6, Out: 16, Kernel: $5 \times 5$, Stride: 1 & ReLU \\
 & MaxPool2d & Kernel: $2 \times 2$, Stride: 2 & - \\
 & Linear & In: 16 $\times H' \times W'$, Out: 120 & ReLU \\
 & Linear & In: 120, Out: 84 & ReLU \\
 & Linear & In: 84, Out: C (Classes) & Softmax / None \\
\midrule
\textbf{ResNet-18} & Conv2d & Input Conv, Kernel $7\times7$, stride 2 & ReLU \\
 & MaxPool2d & Kernel $3\times3$, stride 2 & - \\
 & Residual Blocks & 4 stages of basic blocks (2 convs per block with skip) & ReLU \\
 & AveragePool & Global average pooling & - \\
 & Linear & Fully connected to class outputs & Softmax / None \\
\bottomrule
\end{tabular}
\end{table}

\subsection{Training Setup and Optimization}
% Question 2.3
We train our models using mini-batch Stochastic Gradient Descent (SGD) with momentum. 
\begin{itemize}
    \item \textbf{Optimization Method:} SGD with a momentum factor (typically $0.9$) and weight decay for regularization.
    \item \textbf{Loss Function:} Cross-entropy loss for multi-class classification, or binary cross-entropy (BCE) for the binary classification safety study.
    \item \textbf{Training Hardware and Times:} Training is performed using a hardware accelerator (CUDA/MPS device) to minimize computation time. The total training time and training speeds (epochs per second) are documented in Table~\ref{tab:training_metrics}.
\end{itemize}

\begin{table}[H]
\centering
\caption{Training and Optimization Settings}
\label{tab:training_metrics}
\begin{tabular}{lccccc}
\toprule
Architecture & Batch Size & Optimizer & Learning Rate & Total Epochs & Training Time (s) \\
\midrule
LeNet-5 & - & SGD (momentum) & - & - & - \\
ResNet-18 & - & SGD (momentum) & - & - & - \\
\bottomrule
\end{tabular}
\end{table}

\subsection{Hyperparameter Tuning and Validation Results}
% Question 2.4
Hyperparameter tuning is carried out using a validation set. We evaluate different learning rates, batch sizes, and model depths to avoid overfitting. 
Validation loss and accuracy are monitored at each epoch to implement early stopping.
The validation curves and final test scores are summarized below.

% [Include placeholder for validation/test accuracy comparison plots and tables]

\section{Machine Learning and AI Safety}

\subsection{Real-World Implications of Bias}
% Question 3.1
In real-world applications, machine learning models are often exposed to training datasets that reflect historical inequalities, selection bias, or spurious environmental correlations. For instance:
\begin{itemize}
    \item \textbf{Medical Diagnostics:} If a model is trained on dermatological images where malignant lesions are disproportionately photographed next to a ruler, the model might rely on the presence of the ruler rather than the lesion's features to make predictions.
    \item \textbf{Recruiting Systems:} If historically most successful candidates for a technical role were male, the model might learn to associate gender-specific terms or hobbies with positive outcomes, unfairly penalizing female candidates.
\end{itemize}
These biases degrade model reliability, lead to systemic discrimination, and can have critical safety or legal consequences, especially under regulations like the EU AI Act \cite{besse2022european} and CNIL guidelines \cite{cnil2026guide}.

\subsection{Biased Dataset Construction}
% Question 3.2
To study these vulnerabilities empirically, we construct a biased version of the TILDA dataset for the binary classification task (e.g., cauliflower vs. head cabbage, coded as $y \in \{0, 1\}$). We concatenate a noise-like image $\epsilon \in \mathbb{R}^{N \times N}$ to the original image $x \in \mathbb{R}^{N \times N}$ to build a multi-channel sample $[x, \epsilon]$.

The value of $\epsilon$ is strongly correlated with the true label $y \in \{0, 1\}$ through a bias variable $S$:
\begin{equation}
    S | \{y = k\} \sim \text{Bernoulli}(p_k), \quad k \in \{0, 1\}
\end{equation}
The noise variable $\epsilon$ is then defined as:
\begin{equation}
    \epsilon = \begin{cases} 
        0 & \text{if } S = 0 \\ 
        \mathbf{N}(0, I) & \text{if } S = 1 
    \end{cases}
\end{equation}
By setting $p_0$ and $p_1$, we control the level of spurious correlation. In the extreme case of $p_0 = 0$ and $p_1 = 1$, the noise $\epsilon$ becomes a perfect predictor of $y$, allowing the model to completely ignore the visual features of $x$.

\subsection{Model Bias Evaluation and Experimental Analysis}
% Question 3.3, 3.4, 3.5
We train and evaluate two settings:
\begin{enumerate}
    \item \textbf{Model 1 ($model_1$):} Trained on an unbiased TILDA dataset ($p_0 = 1/2, p_1 = 1/2$).
    \item \textbf{Model 2 ($model_2$):} Trained on a biased version of TILDA ($p_0 = 0, p_1 = 1$).
\end{enumerate}

\subsubsection{Disparate Impact Study}
% Question 3.3
To quantify the bias of the trained models on test data, we compute the Disparate Impact (DI) metric \cite{besse2022european}:
\begin{equation}
    \text{DI} = \frac{\mathbb{P}(\hat{y}(X) = 1 | S = 0)}{\mathbb{P}(\hat{y}(X) = 1 | S = 1)}
\end{equation}
A model is considered unbiased if the DI metric is close to $1.0$. If DI deviates significantly from 1, the model is exploiting the biased noise channel $\epsilon$ to make predictions. We evaluate whether $model_2$ displays strong bias compared to $model_1$.

\subsubsection{Accuracy Comparisons}
% Question 3.4
We evaluate the classification accuracy of both $model_1$ and $model_2$ on the test set. 
Since $model_2$ may learn to ignore $x$ in favor of the easier-to-learn noise signal $\epsilon$, we measure how its accuracy drops when evaluated on unbiased test data or when the correlation between $S$ and $y$ is broken.

\subsubsection{Summary of Results}
% Question 3.5
Table~\ref{tab:bias_evaluation_results} summarizes our findings.

\begin{table}[H]
\centering
\caption{Model Bias and Performance Metrics}
\label{tab:bias_evaluation_results}
\begin{tabular}{lccccc}
\toprule
Model & Training Settings & Train Accuracy & Test Accuracy & DI (Test) \\
\midrule
$model_1$ & Unbiased ($p_0=1/2, p_1=1/2$) & - & - & - \\
$model_2$ & Biased ($p_0=0, p_1=1$) & - & - & - \\
\bottomrule
\end{tabular}
\end{table}

% Discussion and conclusions on learning shortcut/spurious features.

\subsection{Bias Study in the Literature}
% Question 3.6 (bonus)
% Essay addressing: "Est-ce qu'on peut faire confiance à mon modèle ?"
Here, we present a critical discussion on model reliability and bias from the perspective of the EU AI Act compliance and CNIL guides. We discuss how statistical indicators like Disparate Impact are utilized, the limitations of empirical verification, and the steps required to transition towards trustworthy AI systems.

\printbibliography

\end{document}
